\documentclass[10pt,a4paper]{amsart}
\usepackage[margin=19mm]{geometry}
\usepackage{amsmath,amssymb,mathtools}
\usepackage[scr=rsfs,cal=euler]{mathalfa}
\usepackage{xfrac}
\usepackage[unicode,hidelinks,bookmarks=false]{hyperref}
\newcommand{\ie}{i.e.\ }
\newcommand{\cf}{cf.\ }
\newcommand{\resp}{respectively\ }
\newcommand{\Subsection}{\subsection}
\providecommand{\nobreakdash}{\nobreak\hbox{-}}
\setlength{\emergencystretch}{2em}
\multlinegap0pt
\usepackage{bm}
\usepackage{amssymb}
\usepackage{tikz}
\usepackage{enumerate}
\usepackage[shortlabels]{enumitem}
\usepackage{csquotes}
\usepackage{float}
\usepackage{xcolor}
\usepackage{mathtools}
\newtheorem{lemma}{Lemma}
\newtheorem{claim}{Claim}
\newtheorem{poc}{Poc}
\title{A solution to a strengthened conjecture of Bukh, van Hintum and Keevash on additive bases}
\author{Zixiang Xu}
\date{Source: arXiv:2605.10771v1 (CC BY 4.0)}
\begin{document}
\maketitle

\noindent\emph{Source and licence.} A solution to a strengthened conjecture of Bukh, van Hintum and Keevash on additive bases. Version arXiv:2605.10771v1 (CC BY 4.0), distributed by arXiv under the Creative Commons Attribution 4.0 International licence, http://creativecommons.org/licenses/by/4.0/, as recorded in the arXiv licence metadata for this exact version and shown on the abstract page https://arxiv.org/abs/2605.10771v1. Full licence notice: \texttt{licenses/arxiv-2605.10771-CC-BY-4.0.txt}.\par\smallskip

\noindent\textit{Editorial scope.} Counted unit: the article's lemma lem:coset (source0.tex lines 301--314). The article's complete original proof of it is delivered with it (source0.tex lines 316--465, one \texttt{proof} environment of 150 lines). No mathematics is written, repaired or re-derived: the statement and the proof are the article's own lines, in the article's own notation, and no retained line is substituted or re-flowed.  No further source-local context is needed: every cross-reference of every retained line resolves inside the two delivered spans.  Nothing was omitted from any retained span, so the package carries no omission record.  No retained line contains a citation, so the wrapper carries no bibliography and no bibliography entry.  No retained line contains a figure environment, an image-inclusion command or drawing code, so no image is delivered and the package declares no asset dependency.  Editorial formatting only, in addition: an AMS \texttt{amsart} preamble that restates the article's own declarations and macros used by the retained lines, namely the environments \texttt{lemma}, \texttt{claim}, \texttt{poc} on separate counters, together with the standard packages those lines need.  The retained environments print in this document as Lemma 1, Claim 1, Poc 1 in order of appearance, whereas the article prints the same items with the numbering of its own full document.  The complete native source of the article as distributed in the arXiv e-print for this exact version is bundled unchanged as \texttt{sources/200317/source0.tex}.
\begin{lemma}\label{lem:coset}
Let \(n\ge 2\) and \(\Gamma\) be an abelian group, and let \(K\le \Gamma\) be a subgroup identified with \(\mathbb{F}_2^n\). Write
\(
    D_n=\{\boldsymbol{e}_i+\boldsymbol{e}_j:1\le i<j\le n\}\subseteq K.
\)
Let \(u\) be an integer with \(0\le u\le n\). Suppose that \(X\subseteq C\subseteq \Gamma\), \(|X|\le n-u\), and
\(
    D_n\subseteq X-C.
\)
Then
\(
    |C|\ge n+\binom{u}{2}.
\)
\end{lemma}

\begin{proof}[Proof of Lemma~\ref{lem:coset}]
First suppose that \(X,C\subseteq K=\mathbb{F}_2^n\). Since \(n\ge 2\), the set \(D_n\) is nonempty, and hence the containment \(D_n\subseteq X-C\) implies that \(X\) is nonempty. Choose \(\boldsymbol{x}_0\in X\) and replace \(X\) and \(C\) by \(X-\boldsymbol{x}_0\) and \(C-\boldsymbol{x}_0\), respectively. This preserves both the containment \(X\subseteq C\) and the difference set \(X-C\), so we may assume that \(\boldsymbol{0}\in X\).

Put \(m=|X|\), and let
\(
    W=\operatorname{span}_{\mathbb{F}_2}(X).
\)
Since \(\boldsymbol{0}\in X\), the space \(W\) is generated by the at most \(m-1\) nonzero elements of \(X\), and hence
\(
    \dim W\le m-1.
\)
Let
\(
    \rho:\mathbb{F}_2^n\to \mathbb{F}_2^n/W
\)
be the quotient map, and put \(s=\dim(\mathbb{F}_2^n/W)\). Then
\(
    s=n-\dim W\ge n-m+1.
\)
The point of passing to the quotient is that all elements of \(X\) collapse to the zero coset.  We shall use the following consequence of the covering assumption \(D_n\subseteq X-C\).
\begin{claim}\label{claim:containmentRho}
    \(
        \rho(D_n)\subseteq \rho(C).
    \)
\end{claim}
\begin{poc}
    Let \(\boldsymbol{d}\in D_n\).  By the assumption \(D_n\subseteq X-C\), we may write
    \(
        \boldsymbol{d}=\boldsymbol{x}-\boldsymbol{c}
    \)
    for some \(\boldsymbol{x}\in X\) and \(\boldsymbol{c}\in C\).  Since \(X\subseteq W\), we have \(\rho(\boldsymbol{x})=\boldsymbol{0}\), and therefore
    \[
        \rho(\boldsymbol{d})
        =
        \rho(\boldsymbol{x})-\rho(\boldsymbol{c})
        =
        -\rho(\boldsymbol{c})
        =
        \rho(\boldsymbol{c}),
    \]
    where the last equality uses that \(\mathbb{F}_2^n/W\) has exponent \(2\).  Thus \(\rho(\boldsymbol{d})\in \rho(C)\), as claimed.
\end{poc}

We now count the elements of \(C\) according to the \(W\)-cosets they occupy.  Since \(X\subseteq C\cap W\), the zero coset of \(W\) already contains at least the \(m\) elements of \(X\).  Moreover, by Claim~\ref{claim:containmentRho}, every nonzero element of \(\rho(D_n)\) corresponds to a nonzero \(W\)-coset that meets \(C\).  Hence
\[
    |C|\ge m+|\rho(D_n)\setminus\{\boldsymbol{0}\}|.
\]

Since the vectors
\(
    \rho(\boldsymbol{e}_1),\dots,\rho(\boldsymbol{e}_n)
\)
span \(\mathbb{F}_2^n/W\), we may choose indices
\(
    1\le i_1<\cdots<i_s\le n
\)
such that
\(  \rho(\boldsymbol{e}_{i_1}),\dots,\rho(\boldsymbol{e}_{i_s})
\)
form a basis of \(\mathbb{F}_2^n/W\).  The pairwise sums of these \(s\) basis vectors are nonzero and pairwise distinct. For every \(1\le a<b\le s\),
\[
    \rho(\boldsymbol{e}_{i_a})+\rho(\boldsymbol{e}_{i_b})
    =
    \rho(\boldsymbol{e}_{i_a}+\boldsymbol{e}_{i_b})
    \in \rho(D_n).
\]
Hence
\(
    |\rho(D_n)\setminus\{\boldsymbol{0}\}|
    \ge \binom{s}{2}.
\)
Moreover, since \(s=n-\dim W\ge n-m+1\) and \(m\le n-u\), we have
\[
    |C|\ge m+\binom{s}{2}\ge m+\binom{n-m+1}{2}= n+\binom{n-m}{2}\ge n+\binom{u}{2}.
\]
It remains to reduce the general case \(X\subseteq C\subseteq \Gamma\) to the case already proved, in which both sets lie inside \(K=\mathbb{F}_2^n\). The key observation is that, since every element of \(D_n\) lies in \(K\), any representation of an element of \(D_n\) as a difference \(\boldsymbol{x}-\boldsymbol{c}\), with \(\boldsymbol{x}\in X\) and \(\boldsymbol{c}\in C\), must use two elements from the same coset of \(K\).  We therefore fold each \(K\)-coset meeting \(X\) back onto \(K\).

For each coset \(\alpha+K\) meeting \(X\), choose an element
\(
    \boldsymbol{x}_\alpha\in X\cap(\alpha+K),
\)
and define
\[
    X_\alpha'
    =
    (X\cap(\alpha+K))-\boldsymbol{x}_\alpha
    \subseteq K,
    \qquad
    C_\alpha'
    =
    (C\cap(\alpha+K))-\boldsymbol{x}_\alpha
    \subseteq K.
\]
Let
\(
    X'=\bigcup_\alpha X_\alpha'\) and
    \(
    C'=\bigcup_\alpha C_\alpha',
\)
where the unions range over all cosets of \(K\) meeting \(X\).  Since \(X\subseteq C\), we have
\(
    X'\subseteq C'\subseteq K.
\)
Moreover,
\(
    |X'|\le |X|\) and
    \(
    |C'|\le |C|,
\)
because translation preserves cardinalities within each coset, while sets coming from distinct cosets may overlap after being translated into \(K\). We then claim that the covering property of \(D_n\) is preserved under this folding operation.
\begin{claim}\label{claim:DXC}
    \(D_n\subseteq X'-C'.\)
\end{claim}
\begin{poc}
    Let \(\boldsymbol{d}\in D_n\).  Since \(D_n\subseteq X-C\), we write
    \(
        \boldsymbol{d}=\boldsymbol{x}-\boldsymbol{c}
    \)
    for some \(\boldsymbol{x}\in X\) and \(\boldsymbol{c}\in C\).  As \(\boldsymbol{d}\in K\), the elements \(\boldsymbol{x}\) and \(\boldsymbol{c}\) lie in the same coset \(\alpha+K\).  Therefore
    \(
        \boldsymbol{x}-\boldsymbol{x}_\alpha\in X_\alpha'\) and
        \(
        \boldsymbol{c}-\boldsymbol{x}_\alpha\in C_\alpha',
    \)
    which further implies that
    \(
        \boldsymbol{d}
        =
        (\boldsymbol{x}-\boldsymbol{x}_\alpha)
        -
        (\boldsymbol{c}-\boldsymbol{x}_\alpha)
        \in X'-C'.
    \) This finishes the proof.
\end{poc}
By construction and Claim~\ref{claim:DXC}, the pair \(X',C'\) lies inside \(K\) and satisfies
\(
    X'\subseteq C',\) \(
    |X'|\le |X|\le n-u\) and
    \(D_n\subseteq X'-C'.
\)
Therefore the special case already proved for subsets of \(K=\mathbb{F}_2^n\) applies to \(X'\) and \(C'\), and yields
\(
    |C'|\ge n+\binom{u}{2}.
\)
Finally, the folding operation used to construct \(C'\) can only decrease cardinality, so 
\(
    |C|\ge |C'|\ge n+\binom{u}{2}.
\)
This completes the proof.
\end{proof}

\end{document}
